% Anwendungen des allgemeinen Elternsystem-Kriteriums aus Band 26.
\providecommand{\BXXVIIFiniteRooted}[3]{%
  \mathsf{EndWBaum}(#1,#2,#3)}
\providecommand{\BXXVIIParentSystem}[4]{%
  \mathsf{ElterSys}(#1,#2,#3,#4)}
\providecommand{\BXXVIISymmetricEdges}[2]{%
  \operatorname{sym}_{#1}(#2)}
\providecommand{\BXXVIIForwardEdges}[1]{%
  \operatorname{vor}(#1)}


\section{Die gerichtete Darstellung der Positionsbäume}

Beim Positionsbaum ist die Höhenmarke bereits vorhanden: Es ist die
Wortlänge der Adresse. Das Anhängen eines Bits ist eine gerichtete
Elternkante; erst durch das Vergessen ihrer Richtung erhält man die
bisher verwendete symmetrische Adressrelation.

\noindent\begin{minipage}{\linewidth}
\FormulaDefDeltaK[Vorwärtskanten einer Positionsmenge]{%
  \BXXVIIForwardEdges P\coloneqq
  \{(p,q)\in P\times P\mid
      \exists i\in\{0,1\}\,q=\WordConcat p{\LetterWord i}\}
}{B27ForwardAddressEdgesDef}{%
  \DeltaRow{Mengen}{P\dsep p\dsep q\dsep i}
  \DeltaPrem{Positionsmenge}{\GoodAddressSet P}
  \DeltaRow{Neue Menge}{\BXXVIIForwardEdges P}
}
\end{minipage}\par

\noindent\begin{minipage}{\linewidth}
\FormulaThmDeltaK[Gerichtete und ungerichtete Positionskanten]{%
  \begin{aligned}[t]
    \GoodAddressSet P\vdash{}&
      \BXXVIISymmetricEdges P{\BXXVIIForwardEdges P}=\AddressEdges P
  \end{aligned}
}{B27ForwardAddressSymmetrization}{%
  \DeltaRow{Mengen}{P\dsep p\dsep q\dsep i}
}
\end{minipage}\par
\noindent\textit{Beweis.}
Für \(p,q\in P\) bedeutet die Mitgliedschaft links, dass entweder
\(q=p\frown\LetterWord i\) für ein Bit \(i\) gilt oder
\(p=q\frown\LetterWord i\) für ein Bit \(i\) gilt.
Durch Fallunterscheidung ist dies gleichbedeutend mit der Existenz
eines Bits, für das eine dieser beiden Gleichungen gilt. Nach
\FormulaRefAuto{B27AddressEdgeCriterion}[thm] ist dies genau die
Mitgliedschaft in \(\AddressEdges P\). Außerhalb von \(P\times P\)
enthält keine der beiden Relationen ein Paar. Die Mengenextensionalität
liefert die Gleichheit.

\noindent\begin{minipage}{\linewidth}
\FormulaThmDeltaK[Positionsmengen erfüllen das allgemeine Elternsystem]{%
  \begin{aligned}[t]
    &\GoodAddressSet P\dsep h\colon P\to\mathbb N\dsep
      \forall p\in P\,(h(p)=\WordLength p)\vdash\\[-2pt]
    &\BXXVIIParentSystem P{\BXXVIIForwardEdges P}
      {\EmptyWord{\{0,1\}}}h
  \end{aligned}
}{B27AddressRankedParentSystem}{%
  \DeltaRow{Mengen}{P\dsep h\dsep p\dsep q\dsep i}
}
\end{minipage}\par
\noindent\textit{Beweis.}
\FormulaRefAuto{GoodBinaryAddressSetCriterion}[thm] liefert die
Endlichkeit von \(P\), die Wurzel und die Existenz eines unmittelbaren
Präfixes jeder Nichtwurzeladresse. Die Vorwärtskanten liegen nach ihrer
Definition in \(P\times P\). Aus
\FormulaRefAuto{EmptyWordLengthZero}[thm] folgt die Höhenmarke null
an der Wurzel. Jeder Vorgänger einer Nichtwurzeladresse \(q\) ist
\(\WordRumpf{\{0,1\}}q\), denn
\FormulaRefAuto{WordRumpfAppend}[thm] entfernt das angehängte Bit.
Der Vorgänger ist somit eindeutig. Schließlich liefert
\FormulaRefAuto{WordAppendLetterLength}[thm] für jede Vorwärtskante
\(h(q)=h(p)+1\), also das verlangte strenge Wachstum. Dies sind
sämtliche Bedingungen von
\FormulaRefAuto{B27RankedParentSystemDef}[def].

Eine solche Funktion \(h\) existiert: Jedem \(p\in P\) wird die
eindeutig bestimmte natürliche Zahl \(\WordLength p\) zugeordnet;
\FormulaRefAuto{B27AddressWordTyping}[thm] und
\FormulaRefAuto{FiniteWordLengthNatural}[thm] rechtfertigen die
Typisierung. Das Funktionsdefinitionsprinzip bildet daraus ihren Graphen.
Damit ist das allgemeine Einführungskriterium auf jede der zuvor
konstruierten Positionsmengen anwendbar.

Die volle binäre Verzweigung ist eine zusätzliche Bedingung. Im
allgemeinen Elternsystem sind beliebig viele Kinder möglich. Bei
\(\GoodAddressSet P\) kommen als Kinder einer Adresse nur ihre beiden
Erweiterungen um \(0\) und \(1\) infrage. Der Vollheitsabschluss
fordert, dass beide vorhanden sind oder beide fehlen. Die Funktionen
\(\AddressLeft P\) und \(\AddressRight P\) unterscheiden diese Kinder.
Genau dies hat
\FormulaRefAuto{GoodBinaryAddressSetOrderedFullBinaryTree}[thm] als
Struktur aus Band~26 nachgewiesen. Eine spätere Syntax darf daher
den allgemeinen endlichen Wurzelbaum verwenden und ihre jeweiligen
Konstruktoren durch zusätzliche Beschriftungs- und Verzweigungsregeln
kennzeichnen.
